\documentclass[a4paper,oneside,12pt]{amsart}

%\usepackage[spanish, activeacute]{babel}
\usepackage[utf8]{inputenc}
%\usepackage[latin1]{inputenc}
\usepackage{latexsym}
\usepackage{multicol}
\usepackage{enumerate}
\usepackage{booktabs}
\usepackage[heightrounded]{geometry}
\geometry{top=2cm,bottom=2cm,left=2cm,right=2cm}
\usepackage{graphicx}

\DeclareMathOperator{\cond}{cond}
\DeclareMathOperator{\Id}{Id}
\DeclareMathOperator{\re}{Re}
%\DeclareMathOperator{\im}{Im}
\newcommand{\I}{\mathbb{I}}
\newcommand{\R}{\mathbb{R}}
\newcommand{\Q}{\mathbb{Q}}
\newcommand{\C}{\mathbb{C}}
\newcommand{\Z}{\mathbb{Z}}
\newcommand{\N}{\mathbb{N}}
\DeclareMathOperator{\im}{Im}

\newcommand{\nombremateria}
{Investigaci\'on Operativa}
\newcommand{\nombrecuatrimestre}
    {Segundo cuatrimestre de 2018}
\newcommand{\numeroparcial}
    { Segundo Parcial }
\newcommand{\fecha}
    {20 de Octubre de 2020}
\newcommand{\numerotema}
    {\bf 1}

%\oddsidemargin -1cm \evensidemargin -1cm \textwidth 17.5cm
%topmargin 1.2cm \textheight 25cm
%\setlength{\textheight}{25cm}

\def \ds{\displaystyle}
\newtheorem{quest}{Ejercicio}
\thispagestyle{empty}

\begin{document}

\hfill \hfill
\begin{tabular}[c]{|c|c|c|c|c|}
\hline 
1 & 2 & 3 & 4 & Calificaci\'on \\ \hline \hline
\hspace{1 cm} & \hspace{1 cm} & \hspace{1 cm} & \hspace{1 cm} & \hspace{1.3 cm} \\
\hspace{1 cm} & \hspace{1  cm} & \hspace{1  cm} & \hspace{1  cm} & \hspace{1.3 cm} \\ \hline

\end{tabular}

\vspace{-1cm}

\noindent Nombre y apellido:

\vspace{0.2cm}

\noindent Número de libreta:

\noindent \hrulefill 

\smallskip
\renewcommand{\baselinestretch}{1.1}

\begin{center}
	\large \textbf{\nombremateria} 
	
	\numeroparcial -- \fecha
\end{center}

% \vspace{.1cm}
\renewcommand{\theenumi}{\textbf{\arabic{enumi}}}

%%% EMPIEZA EL EJERCICIO 1 %%%

\begin{quest} Resolver el siguiente problema de Programación Lineal utilizando el Método de Dos Fases de SIMPLEX.
    \[\begin{array}{rrrrrrrr}
        \min & 5x_1 & - & x_2 & - & 2x_3  \\
        \text{s.a:} & x_1 & + & x_2 & - & x_3 & = & 6 \\
          & x_1 & + & 2x_2 & - & x_3 & \geq & 4 \\
          & -x_1  & + & x_2 & + & 2x_3 & \leq & 8 \\
         & \multicolumn{7}{l}{x_1\leq 0 \;\;\;x_2,x_3\geq 0}
    \end{array}
    \]

\end{quest}

%%% TERMINA EL EJERCICIO 1 %%%

%%% EMPIEZA EL EJERCICIO 2 %%%
\begin{quest}
Cada uno de los siguientes diccionarios corresponden a soluciones básicas de distintos problemas de Programación Lineal con objetivo de maximizar. 
% En cada caso, seleccionar qué información puede extraerse sin realizar ningún cálculo y justificar la elección.
% En cada caso se presentan cuatro afirmaciones sobre la información que brinda ese diccionario. Elegir cuál o cuáles de esas afirmaciones son verdaderas y justificar la elección (no es necesario justificar por qué no se eligieron las demás). En todos los casos puede considerarse que se utiliza el criterio usual para la selección de la variable que entrará a la base en la siguiente iteración.
En cada caso se presentan cuatro afirmaciones sobre la información que brinda ese diccionario. Decidir y justificar si cada afirmación es verdadera o falsa. En todos los casos puede considerarse que se utiliza el criterio usual para la selección de la variable que entrará a la base en la siguiente iteración.
\begin{itemize}
    \item[(1)]
            \[\begin{array}{rrrrrrrrr}
                x_1 & = & 2 & + & 5x_4 & - & 2x_3 & + & x_6 \\
                x_2&=&&&x_4&+&x_3&-&x_6\\ 
                x_5&=&3&&&-&\frac{1}{2}x_3&&\\ \hline
                z&=&4&+&\frac{1}{2}x_4&&&+& x_6
            \end{array}\]
            \begin{enumerate}
                \item[(a)] corresponde a una solución básica óptima
                \item[(b)] la próxima iteración de SIMPLEX es degenerada
                \item[(c)] corresponde a una solución básica degenerada
                \item[(d)] el problema es no acotado
            \end{enumerate}
    % \item[(2)]
    %         \[
    %         \begin{array}{rrrrrrr}
    %             x_3&=&&&2x_2&+&x_5\\
    %             x_1&=&2&+&x_2&-&x_5\\
    %             x_4&=&1&-&x_2&-&2x_5\\ \hline
    %             z&=&&&x_2&+&\frac{1}{2}x_5
    %         \end{array}
    %         \]
    %         \begin{enumerate}
    %             \item[(a)] la próxima iteración de SIMPLEX es degenerada
    %             \item[(b)] el problema tiene infinitas soluciones óptimas
    %             \item[(c)] el problema es no acotado
    %             \item[(d)] en la próxima iteración de SIMPLEX, $x_4$ sale de la base
    %         \end{enumerate}
    \item[(2)]
            \[
            \begin{array}{rrrrrrr}
                x_5&=&\frac{1}{2}&-&x_2&+&x_4\\
                x_1&=&2&+&\frac{1}{2}x_2&-&\frac{3}{2}x_4\\
                x_3&=&1&-&\frac{1}{2}x_2&&\\ \hline
                z&=&11&&&&\\
            \end{array}
            \]
            \begin{enumerate}
                \item[(a)] la próxima iteración de SIMPLEX es degenerada
                \item[(b)] corresponde a una solución básica óptima
                \item[(c)] existen infinitas soluciones óptimas
                \item[(d)] el problema dual es infactible
            \end{enumerate}

\end{itemize}
\end{quest}
%%% TERMINA EL EJERCICIO 2 %%%

%%% EMPIEZA EL EJERCICIO 3 %%%


\begin{quest}
Dado el siguiente problema de Programación Lineal:
\[
\begin{array}{rrrrrrrrrr}
    \max & -(2+\alpha)x_1 & + & (1+\alpha) x_2 & - & 4x_3 & + & 9x_4 & & \\
    \text{s.a:} & 2x_1 &  &  & + & x_3& +& x_4 & = & 2 + \beta \\
     & -x_1& + &x_2 &+ &2x_3& + &\frac{1}{2}x_4& \leq & 5 + \gamma \\
     & \multicolumn{9}{l}{x_1,x_2,x_3\geq 0\, , x_4\leq 0}
\end{array}
\]
\begin{itemize}
    \item[(a)] Para el caso $\alpha=\beta=\gamma=0$, hallar una solución óptima utilizando el problema dual.
    \item[(b)] Hallar todos los $\alpha\in\mathbb{R}$ para los cuales la solución obtenida en el ítem (a) sigue siendo óptima.
    \item[(c)] Suponiendo que $\alpha$ cumple con las condiciones del ítem (b), hallar condiciones para que la solución obtenida en el ítem (a) siga siendo óptima cuando se mueven a la vez $\beta$ y $\gamma$.
\end{itemize}
\end{quest}

%%% TERMINA EL EJERCICIO 3 %%%

%%% EMPIEZA EL EJERCICIO 4 %%%


\begin{quest} 
Utilizar el algoritmo de Branch \& Bound para resolver el siguiente problema de programaci\'on lineal entera, detallando cada divisi\'on en subproblemas durante el algoritmo mediante el esquema de \'arbol visto en clase. Comenzar ramificando en $x_2$ y resolver gráficamente la relajación lineal de cada subproblema. 

\[\begin{array}{rrrrrrrrc}
\max & 2x_1 & + & x_2 &  &   \\
\text{s.a:}	 & 5x_1 & - & 5x_2 & \geq & 1  \\
			 & -x_1 & + & 2x_2 & \geq & -2 \\
			 &2x_1 & +  & 3x_2 & \leq & 45 \\
			 &	2x_1    & +  & x_2 & \geq & 2 \\
			 & \multicolumn{3}{r}{x_1,x_2} & \in & \mathbb{Z}
\end{array}
\]



\end{quest}
%%% TERMINA EL EJERCICIO 4 %%%

%%% TERMINA EL PARCIAL %%%



\end{document}

